\documentclass[../../main.tex]{subfiles}
\begin{document}
\Needspace{12\baselineskip}
% SOURCE p=192; printed=177; object=ch05-hilbert-exercises-title; region=middle
\subsection*{习题}
% SOURCE p=192; printed=177; object=ch05-ex54-ell2-completeness; number=54; region=middle
\begin{exercise}\label{exercise:ch05-ell2-completeness}
对任意非空集合 $A$，$\ell^2(A)$ 完备。
\end{exercise}
% SUPPLEMENT: id=v14-ch05-ex54-solution; parent=ch05-ex54-ell2-completeness; kind=solution+motivation+hypotheses; origin=independent-derivation
\begin{solution}
\textbf{逐坐标完整补写（重构新增）。}设 $(f_n)$ 为 $\ell^2(A)$ 中的 Cauchy 列。对每个 $\alpha$，$|f_n(\alpha)-f_m(\alpha)|\leqslant\|f_n-f_m\|_2$，故标量列有极限 $f(\alpha)$。Cauchy 列有界，记 $M=\sup_n\|f_n\|_2<\infty$。对任意有限 $F\subseteq A$，
\begin{align*}
\sum_{\alpha\in F}|f(\alpha)|^2
=\lim_n\sum_{\alpha\in F}|f_n(\alpha)|^2\leqslant M^2.
\end{align*}
对有限 $F$ 取上确界得 $f\in\ell^2(A)$。给 $\varepsilon>0$，选 $N$ 使 $n,m\geqslant N$ 时 $\|f_n-f_m\|_2<\varepsilon/2$。固定 $n\geqslant N$，对有限 $F$ 令 $m\to\infty$，得 $\sum_F|f_n-f|^2\leqslant\varepsilon^2/4$，再取上确界得 $\|f_n-f\|_2\leqslant\varepsilon/2<\varepsilon$。故整列收敛。
\end{solution}
\begin{remark}
\textbf{动机与条件说明（重构新增）。}先从每个坐标的标量完备性获得候选极限，再用同一个有限和上界控制所有坐标。不能只凭逐坐标收敛断言范数收敛；Cauchy 性的统一误差是关键。$A$ 可以不可数，证明从未枚举它。空 $A$ 对应零空间，也成立。
\end{remark}
% END SUPPLEMENT: id=v14-ch05-ex54-solution
% SOURCE p=192; printed=177; object=ch05-ex55-polarization-unitary; number=55; region=middle
\begin{exercise}\label{exercise:ch05-polarization-unitary}
设 $\mathscr H$ 为 Hilbert 空间。
\begin{enumerate}
\item\label{item:ch05-polarization-identity} （极化恒等式）对任意 $x,y\in\mathscr H$，
\begin{align*}
\langle x,y\rangle=\frac14\bigl(\|x+y\|^2-\|x-y\|^2
+\mathrm i\|x+\mathrm iy\|^2-\mathrm i\|x-\mathrm iy\|^2\bigr).
\end{align*}
这里不需要完备性。
\item\label{item:ch05-unitary-isometric-surjective} 若 $\mathscr H'$ 为另一个 Hilbert 空间，则从 $\mathscr H$ 到 $\mathscr H'$ 的线性映射为酉映射，当且仅当它等距且满射。
\end{enumerate}
\end{exercise}
% SUPPLEMENT: id=v14-ch05-ex55-solution; parent=ch05-ex55-polarization-unitary; kind=solution+motivation+hypotheses; origin=independent-derivation
\begin{solution}
\textbf{两问补写（重构新增）。}
\begin{enumerate}
\item 记 $c=\langle x,y\rangle$。内积展开给 $\|x+y\|^2-\|x-y\|^2=4\operatorname{Re}c$。第二变量共轭线性给 $\langle x,\mathrm iy\rangle=-\mathrm ic$，因此 $\|x+\mathrm iy\|^2-\|x-\mathrm iy\|^2=4\operatorname{Im}c$。前一个差加上后一个差的 $\mathrm i$ 倍，再除以4，恰为 $c$。
\item 酉映射满射且保持平方范数，故等距。反之，若 $U$ 线性、等距且满射，则 $\|Ux\|=\|x\|$ 使其单射，因此可逆。将第一问的恒等式应用于 $Ux,Uy$，用 $U(x\pm y)=Ux\pm Uy$ 及 $U(x\pm\mathrm iy)=Ux\pm\mathrm iUy$，每个范数都变回原空间中的范数，得到 $\langle Ux,Uy\rangle=\langle x,y\rangle$，故酉。
\end{enumerate}
\end{solution}
\begin{remark}
\textbf{动机与逐条件说明（重构新增）。}范数的四个组合恢复复内积的实部和虚部，因此线性保范数映射保内积。线性条件既保证加减相容，也保证乘 $\mathrm i$ 相容；一般度量等距映射不能替代。满射用于成为两个全空间之间的同构，单射则由等距自动得到。两问的这些计算均不需要完备性。实内积空间使用 $\langle x,y\rangle=(\|x+y\|^2-\|x-y\|^2)/4$，不含虚数项。
\end{remark}
% END SUPPLEMENT: id=v14-ch05-ex55-solution
% SOURCE p=192; printed=177; object=ch05-ex56-double-orthogonal; number=56; region=middle
\begin{exercise}\label{exercise:ch05-double-orthogonal}
若 $E$ 为 Hilbert 空间 $\mathscr H$ 的子集，则 $(E^\perp)^\perp$ 是 $\mathscr H$ 中包含 $E$ 的最小闭子空间。
\end{exercise}
% SUPPLEMENT: id=v14-ch05-ex56-solution; parent=ch05-ex56-double-orthogonal; kind=solution+motivation+hypotheses; origin=independent-derivation
\begin{solution}
\textbf{补写解答（重构新增）。}令 $M=\overline{\operatorname{span}E}$。共轭对称性和线性、连续性给 $E^\perp=M^\perp$：消去 $E$ 等价于消去其有限张成空间，再等价于消去其闭包。显然 $M\subseteq(M^\perp)^\perp$。反过来，取 $x\in(M^\perp)^\perp$，由正交分解写成 $x=m+z$，其中 $m\in M,z\in M^\perp$。因为 $\langle x,z\rangle=0$ 且 $\langle m,z\rangle=0$，得到 $\|z\|^2=0$，所以 $x=m\in M$。故 $(E^\perp)^\perp=M$。任何包含 $E$ 的闭子空间都包含其有限线性组合及闭包，验证最小性。
\end{solution}
\begin{remark}
\textbf{动机与条件说明（重构新增）。}第一次取正交补留下所有测量上不可见的方向，第二次便恢复 $E$ 能生成的闭线性空间，而不只是集合 $E$。$E$ 无须线性或闭；空集时双正交补为零空间。Hilbert 完备性在闭子空间的正交分解中使用，不能把双正交补随意解释成未闭的代数张成空间。
\end{remark}
% END SUPPLEMENT: id=v14-ch05-ex56-solution
% SOURCE p=192; printed=177; object=ch05-ex57-hilbert-adjoint; number=57; region=middle-lower
\begin{exercise}\label{exercise:ch05-hilbert-adjoint}
设 $\mathscr H$ 为 Hilbert 空间，$T\in L(\mathscr H,\mathscr H)$。
\begin{enumerate}
\item\label{item:ch05-hilbert-adjoint-existence}
\begin{definition*}[Hilbert 空间的伴随]\phantomsection\label{definition:ch05-hilbert-adjoint}
存在唯一的 $T^*\in L(\mathscr H,\mathscr H)$，使 $\langle Tx,y\rangle=\langle x,T^*y\rangle$ 对所有 $x,y\in\mathscr H$ 成立，称为 $T$ 的\textbf{伴随}。
\end{definition*}
对照\cref{exercise:ch05-adjoint}。有 $T^*=V^{-1}T^\dagger V$，其中 $V$ 是\cref{theorem:ch05-riesz-hilbert}中从 $\mathscr H$ 到 $\mathscr H^*$ 的共轭线性同构，$(Vy)(x)=\langle x,y\rangle$。
\item\label{item:ch05-hilbert-adjoint-identities} $\|T^*\|=\|T\|$，$\|T^*T\|=\|T\|^2$，$(aS+bT)^*=\overline aS^*+\overline bT^*$，$(ST)^*=T^*S^*$，且 $T^{**}=T$。
\item\label{item:ch05-hilbert-adjoint-range-kernel} 以 $\mathscr R$ 和 $\mathscr N$ 表示值域与零空间，则 $\mathscr R(T)^\perp=\mathscr N(T^*)$，且 $\mathscr N(T)^\perp=\overline{\mathscr R(T^*)}$。
\item\label{item:ch05-unitary-adjoint-inverse} $T$ 为酉算子，当且仅当 $T$ 可逆且 $T^{-1}=T^*$。
\end{enumerate}
\end{exercise}
% SUPPLEMENT: id=v14-ch05-ex57-solution; parent=ch05-ex57-hilbert-adjoint; kind=solution+motivation+hypotheses; origin=independent-derivation
\begin{solution}
\textbf{四问完整补写（重构新增）。}
\begin{enumerate}
\item 固定 $y$，泛函 $x\mapsto\langle Tx,y\rangle$ 线性且模不超过 $\|T\|\|y\|\|x\|$，由 Riesz 表示定理有唯一向量记作 $T^*y$ 表示它，且 $\|T^*y\|\leqslant\|T\|\|y\|$。对任意 $a,b,y,z$，
\begin{align*}
\langle x,T^*(ay+bz)\rangle
=\overline a\langle Tx,y\rangle+\overline b\langle Tx,z\rangle
=\langle x,aT^*y+bT^*z\rangle.
\end{align*}
表示的唯一性给 $T^*$ 线性；前面的界给有界。若另一个算子满足同一恒等式，逐个 $y$ 的唯一性使两者相同。并且 $(T^\dagger Vy)(x)=(Vy)(Tx)=\langle Tx,y\rangle=(VT^*y)(x)$，所以 $VT^*=T^\dagger V$，得到所述公式。
\item 由共轭对称性，$\langle T^*x,y\rangle=\langle x,Ty\rangle$，故 $T^{**}=T$。第一问给 $\|T^*\|\leqslant\|T\|$；应用于 $T^*$ 再给反向不等式。对每个单位向量 $x$，
\begin{align*}
\|Tx\|^2=\langle x,T^*Tx\rangle\leqslant\|T^*T\|.
\end{align*}
取上确界，并结合 $\|T^*T\|\leqslant\|T^*\|\|T\|=\|T\|^2$，得所需等式。零空间时各算子范数都为零，同样成立。最后，对任意 $x,y$，
\begin{align*}
\langle(aS+bT)x,y\rangle
&=\langle x,\overline aS^*y+\overline bT^*y\rangle,\\
\langle STx,y\rangle&=\langle Tx,S^*y\rangle=\langle x,T^*S^*y\rangle.
\end{align*}
由伴随唯一性得两个运算公式。
\item $y\in\mathscr R(T)^\perp$ 等价于 $\langle Tx,y\rangle=0$ 对所有 $x$，等价于 $\langle x,T^*y\rangle=0$ 对所有 $x$，再取 $x=T^*y$ 即等价于 $T^*y=0$。把此结论应用于 $T^*$，得到 $\mathscr R(T^*)^\perp=\mathscr N(T)$；两边取正交补，用上一题及值域的线性性，得第二个等式。
\item 若 $T$ 酉，对任意 $y$ 写 $y=Tz$，则 $\langle Tx,y\rangle=\langle Tx,Tz\rangle=\langle x,z\rangle=\langle x,T^{-1}y\rangle$，故 $T^*=T^{-1}$。反之，若可逆且 $T^*=T^{-1}$，则 $\langle Tx,Ty\rangle=\langle x,T^*Ty\rangle=\langle x,y\rangle$，所以酉。
\end{enumerate}
\end{solution}
\begin{remark}
\textbf{动机与逐条件说明（重构新增）。}伴随把算子从内积的一边移到另一边，Riesz 定理使原本作用于对偶空间的 $T^\dagger$ 回到原空间。两个共轭线性识别与一个线性算子的组合仍线性，而标量移到伴随上须共轭。$T$ 有界保证每个待表示泛函连续；Hilbert 完备性负责表示定理。值域一般不闭，所以第三问的闭包不能省；例如 $\ell^2$ 上 $Tx=(x_n/n)$ 的伴随就是自身，其值域稠密但不等于全空间（$(1/n)$ 不在其值域），零空间却为零。第四问中的可逆性包含满射，单有 $T^*T=I$ 只给等距嵌入，右移算子就是非满的反例。
\end{remark}
% END SUPPLEMENT: id=v14-ch05-ex57-solution
% SOURCE p=192; printed=177; object=ch05-ex58-orthogonal-projection; number=58; region=lower
\begin{exercise}\label{exercise:ch05-orthogonal-projection}
设 $\mathscr M$ 是 Hilbert 空间 $\mathscr H$ 的闭子空间。对 $x\in\mathscr H$，令 $Px$ 是\cref{theorem:ch05-orthogonal-decomposition}中满足 $x-Px\in\mathscr M^\perp$ 的 $\mathscr M$ 中元素。
\begin{enumerate}
\item\label{item:ch05-orthogonal-projection-properties} $P\in L(\mathscr H,\mathscr H)$，并按上一题的记号，有 $P^*=P$，$P^2=P$，$\mathscr R(P)=\mathscr M$，$\mathscr N(P)=\mathscr M^\perp$。
\begin{definition*}[正交投影]\phantomsection\label{definition:ch05-orthogonal-projection}
$P$ 称为到 $\mathscr M$ 上的\textbf{正交投影}。
\end{definition*}
\item\label{item:ch05-self-adjoint-idempotent} 反之，若 $P\in L(\mathscr H,\mathscr H)$ 满足 $P^2=P^*=P$，则 $\mathscr R(P)$ 闭，且 $P$ 为到 $\mathscr R(P)$ 上的正交投影。
\item\label{item:ch05-projection-basis-expansion} 若 $\{u_\alpha\}$ 为 $\mathscr M$ 的标准正交基，则 $Px=\sum_\alpha\langle x,u_\alpha\rangle u_\alpha$。
\end{enumerate}
\end{exercise}
\Needspace{8\baselineskip}
% SUPPLEMENT: id=v14-ch05-ex58-solution; parent=ch05-ex58-orthogonal-projection; kind=solution+motivation+hypotheses; origin=independent-derivation
\begin{solution}
\textbf{三问补写（重构新增）。}
\begin{enumerate}
\item 由分解唯一性，$P(ax+by)=aPx+bPy$，因为剩余项 $a(x-Px)+b(y-Py)$ 仍在 $\mathscr M^\perp$。Pythagoras 给 $\|x\|^2=\|Px\|^2+\|x-Px\|^2$，故 $\|Px\|\leqslant\|x\|$，$P$ 有界。对 $m\in\mathscr M$，其分解为 $m+0$，故 $Pm=m$，从而 $P^2=P$ 且值域正是 $\mathscr M$。$Px=0$ 等价于 $x\in\mathscr M^\perp$，给核。最后，
\begin{align*}
\langle Px,y\rangle=\langle Px,Py\rangle=\langle x,Py\rangle,
\end{align*}
两次等式分别消去正交误差，故 $P^*=P$。
\item 因为 $P^2=P$，其值域等于 $\ker(I-P)$，所以闭。对每个 $x$，$Px$ 在值域中，而对任意 $Pz$，
\begin{align*}
\langle x-Px,Pz\rangle=\langle P(x-Px),z\rangle=0.
\end{align*}
故 $x-Px\in\mathscr R(P)^\perp$，由唯一分解，$P$ 就是相应正交投影。
\item 闭子空间 $\mathscr M$ 本身是 Hilbert 空间，故 $Px=\sum_\alpha\langle Px,u_\alpha\rangle u_\alpha$。因为 $x-Px\perp\mathscr M$，每个系数等于 $\langle x,u_\alpha\rangle$，得到公式，收敛意义与正文正交展开相同。
\end{enumerate}
\end{solution}
\begin{remark}
\textbf{动机与逐条件说明（重构新增）。}正交投影保留子空间分量并丢弃垂直误差，所以长度不增加。子空间闭且环境完备保证每个向量都有这样的分解；$\mathscr M=\{0\}$ 时 $P=0$，否则 $\|P\|=1$，因为它固定某个单位向量。反向中幂等性给代数投影和闭值域，自伴性把这种投影变成正交投影；只有幂等性不能保证正交，例如 $P(x,y)=(x+y,0)$。第三问要求基完整于 $\mathscr M$，无需完整于整个 $\mathscr H$。
\end{remark}
% END SUPPLEMENT: id=v14-ch05-ex58-solution
% SOURCE p=192; printed=177; object=ch05-ex59-minimal-convex-norm; number=59; region=bottom
\begin{exercise}\label{exercise:ch05-minimal-convex-norm}
Hilbert 空间中的每个闭凸集 $K$ 都有唯一的最小范数元素。（若 $0\in K$，结论平凡；否则修改\cref{theorem:ch05-orthogonal-decomposition}的证明。）
\end{exercise}
% SUPPLEMENT: id=v14-ch05-ex59-solution; parent=ch05-ex59-minimal-convex-norm; kind=solution+source-correction+motivation+hypotheses; origin=independent-derivation
\begin{remark}
\textbf{原题缺少非空条件（重构新增）。}空集也闭且凸，却没有任何元素。原题须补 $K\ne\varnothing$；以下证明这一正确版本。
\end{remark}
\begin{solution}
\textbf{完整补写（重构新增）。}设 $K\ne\varnothing$，令 $\delta=\inf_{x\in K}\|x\|$，其为有限非负实数。取 $x_n\in K$ 使 $\|x_n\|\to\delta$。凸性给 $(x_n+x_m)/2\in K$，由平行四边形恒等式，
\begin{align*}
\|x_n-x_m\|^2
=2\|x_n\|^2+2\|x_m\|^2-4\Bigl\|\frac{x_n+x_m}{2}\Bigr\|^2
\leqslant2\|x_n\|^2+2\|x_m\|^2-4\delta^2\longrightarrow0.
\end{align*}
故 $(x_n)$ Cauchy，由空间完备和 $K$ 闭，存在 $x\in K$ 使 $x_n\to x$。范数连续给 $\|x\|=\delta$。若 $y\in K$ 也满足 $\|y\|=\delta$，对 $x,y$ 重复同一恒等式得到 $\|x-y\|^2\leqslant0$，故 $x=y$。
\end{solution}
\begin{remark}
\textbf{动机与逐条件说明（重构新增）。}凸性允许以中点比较两个候选最优点，平行四边形恒等式迫使极小化列互相接近。非空性负责下确界和候选列，闭性保留极限，完备性保证极限存在。删闭性可取 $K=(0,1)\subseteq\mathbb R$，下确界0不取到；删凸性可取 $K=\{-1,1\}$，最小元不唯一。这里不用 $K$ 为线性子空间，也不要求有界；一般最近点问题可对平移集 $K-a$ 应用本结论。上述反例说明条件不能普遍删去，并不声称每个有唯一最小元的集合都必须凸。
\end{remark}
% END SUPPLEMENT: id=v14-ch05-ex59-solution
% SOURCE p=192,193; printed=177,178; object=ch05-ex60-l2-disjoint-decomposition; number=60; region=bottom-top
\begin{exercise}\label{exercise:ch05-l2-disjoint-decomposition}
设 $(X,\mathscr M,\mu)$ 为测度空间。若 $E\in\mathscr M$，把 $L^2(E,\mu)$ 等同于 $L^2(X,\mu)$ 中由在 $E$ 外为零的函数组成的子空间。若 $(E_n)$ 是 $\mathscr M$ 中两两不交的一列集合，且 $X=\bigcup_{n=1}^{\infty}E_n$，则 $\{L^2(E_n,\mu)\}$ 是 $L^2(X,\mu)$ 中一列互相正交的子空间，每个 $f\in L^2(X,\mu)$ 都唯一地写成 $f=\sum_{n=1}^{\infty}f_n$，其中 $f_n\in L^2(E_n,\mu)$，级数按范数收敛。若每个 $L^2(E_n,\mu)$ 可分，则 $L^2(X,\mu)$ 可分。
\end{exercise}
% SUPPLEMENT: id=v14-ch05-ex60-solution; parent=ch05-ex60-l2-disjoint-decomposition; kind=solution+motivation+hypotheses; origin=independent-derivation
\begin{solution}
\textbf{分解、唯一性及可分性补写（重构新增）。}限制测度空间中的函数向 $E$ 外延拓为0，保持 $L^2$ 范数，且与几乎处处等价关系相容，给所述等距识别。对应子空间为有界乘法算子 $f\mapsto\chi_{E^c}f$ 的核，所以闭。不同 $E_n$ 不交使其上支撑的函数乘积为零，故相互正交。

取 $f_n=\chi_{E_n}f$，则 $f_n\in L^2(E_n,\mu)$，且由积分的可数可加性，
\begin{align*}
\Bigl\|f-\sum_{n=1}^Nf_n\Bigr\|_2^2
=\int_{\bigcup_{n>N}E_n}|f|^2\,\mathrm d\mu
=\sum_{n>N}\|f_n\|_2^2\longrightarrow0,
\end{align*}
因为全和为 $\|f\|_2^2<\infty$。若另有范数收敛分解 $f=\sum_ng_n$，对固定 $k$ 应用连续算子 $h\mapsto\chi_{E_k}h$，其作用于足够长的部分和恰为 $g_k$，从而 $g_k=\chi_{E_k}f=f_k$，证明唯一性。

最后选各子空间的可数稠密集 $D_n$ 并添入0。所有有限和 $d_1+\cdots+d_N$（$d_n\in D_n$）构成可数集。先把 $f$ 的分解截到前 $N$ 项，使尾部范数小于 $\varepsilon/2$，再逐项选 $d_n$ 使 $\|f_n-d_n\|_2<\varepsilon/(2N)$，三角不等式给总误差小于 $\varepsilon$。故这些有限和稠密，得到可分性。
\end{solution}
\begin{remark}
\textbf{动机与逐条件说明（重构新增）。}可测分割把积分能量拆成互不干扰的块。可测性保证限制与示性函数运算合法，不交性负责正交，覆盖性保证没有遗漏区域，可数性同时保证级数和可数稠密集构造。允许空块或零测块；不要求各块测度有限，也不要求 $\mu$ 为 $\sigma$ 有限。这里“在 $E$ 外为零”按 $L^2$ 等价类理解，取零延拓代表即可。
\end{remark}
% END SUPPLEMENT: id=v14-ch05-ex60-solution
% SOURCE p=193; printed=178; object=ch05-ex61-product-orthonormal-basis; number=61; region=upper
\begin{exercise}\label{exercise:ch05-product-orthonormal-basis}
设 $(X,\mathscr M,\mu)$ 和 $(Y,\mathscr N,\nu)$ 为 $\sigma$ 有限测度空间，且 $L^2(\mu)$ 与 $L^2(\nu)$ 可分。若 $\{f_m\}$、$\{g_n\}$ 分别为它们的标准正交基，令 $h_{mn}(x,y)=f_m(x)g_n(y)$，则 $\{h_{mn}\}$ 为 $L^2(\mu\times\nu)$ 的标准正交基。
\end{exercise}
% SUPPLEMENT: id=v14-ch05-ex61-solution; parent=ch05-ex61-product-orthonormal-basis; kind=solution+motivation+hypotheses; origin=independent-derivation
\begin{solution}
\textbf{乘积与完全性补写（重构新增）。}取可测代表。由 Tonelli 定理，$\|f_mg_n\|_2^2=\|f_m\|_2^2\|g_n\|_2^2=1$，故 $h_{mn}\in L^2(\mu\times\nu)$。乘积函数的定义与代表选择相容：若某代表在零集上改变，$\sigma$ 有限性将另一空间分解为可数个有限测度集合，可见该零集与另一空间的乘积为乘积零集。Schwarz 不等式使两个 $L^2$ 函数的乘积可积；用 Fubini 分离变量，得
\begin{align*}
\langle h_{mn},h_{pq}\rangle
=\langle f_m,f_p\rangle\langle g_n,g_q\rangle
=\begin{cases}1,&m=p, n=q,\\0,&\text{否则}.
\end{cases}
\end{align*}
这里各因子乘积 $f_m\overline{f_p}$、$g_n\overline{g_q}$ 都绝对可积，因此分离积分合法。

设 $h\in L^2(\mu\times\nu)$ 与所有 $h_{mn}$ 正交。由\cref{theorem:ch02-fubini-tonelli}，几乎每个 $y$ 的截面 $h^y$ 属于 $L^2(\mu)$，且 $\int\|h^y\|_2^2\,\mathrm d\nu(y)=\|h\|_2^2$。固定 $m$，令
\begin{align*}
G_m(y)=\int h(x,y)\overline{f_m(x)}\,\mathrm d\mu(x)
\end{align*}
于截面平方可积处，在其余可测零集上取0。参数积分的可测性由 Tonelli 应用于被积函数实、虚部的正负部分得到；在有效截面上它们都有限，因为 Schwarz 给绝对可积。又有 $|G_m(y)|^2\leqslant\|h^y\|_2^2$，所以 $G_m\in L^2(\nu)$。对每个 $n$，函数 $h\overline{h_{mn}}$ 由 Schwarz 在乘积空间中绝对可积，故 Fubini 给
\begin{align*}
\langle G_m,g_n\rangle=\langle h,h_{mn}\rangle=0.
\end{align*}
$\{g_n\}$ 完全，故 $G_m=0$ 几乎处处。$m$ 的指标集至多可数，取这些例外零集的可数并后，几乎每个 $y$ 同时满足 $\langle h^y,f_m\rangle=0$ 对所有 $m$。由 $\{f_m\}$ 完全得 $h^y=0$ 于 $L^2(\mu)$，再用 Tonelli 得 $\|h\|_2^2=0$，故 $h=0$。标准正交族完全，因此为基。

若某一基为空，相应因子 $L^2$ 为零空间；截面论证直接给乘积 $L^2$ 也为零空间，空乘积族正是其基。有限基的论证同样有效。
\end{solution}
\begin{remark}
\textbf{动机与逐条件说明（重构新增）。}各变量的正交坐标相乘，产生乘积空间的坐标。$\sigma$ 有限性使本书的乘积测度、截面和 Fubini--Tonelli 结论适用；本解不把它宣称为所有可能推广中的必要条件。可分性允许选可数基，因而可把各个“几乎处处”的结论合并在同一个满测集上；不能把不可数个零测集的并当成零测集。两套基的完全性分别消去 $Y$ 和 $X$ 方向，只有标准正交而不完全不足以推出乘积族完整。
\end{remark}
% END SUPPLEMENT: id=v14-ch05-ex61-solution
% SOURCE p=193; printed=178; object=ch05-ex62-lebesgue-l2-separable; number=62; region=upper-middle
\begin{exercise}\label{exercise:ch05-lebesgue-l2-separable}
本题中定义 $L^2$ 空间的测度均为 Lebesgue 测度。
\begin{enumerate}
\item\label{item:ch05-continuous-l2-dense} $C([0,1])$ 在 $L^2([0,1])$ 中稠密。（修改\cref{theorem:ch02-lone-density}的证明。）
\item\label{item:ch05-polynomial-l2-dense} 多项式集合在 $L^2([0,1])$ 中稠密。
\item\label{item:ch05-interval-l2-separable} $L^2([0,1])$ 可分。
\item\label{item:ch05-line-l2-separable} $L^2(\mathbb R)$ 可分。（使用\cref{exercise:ch05-l2-disjoint-decomposition}。）
\item\label{item:ch05-euclidean-l2-separable} $L^2(\mathbb R^n)$ 可分。（使用\cref{exercise:ch05-product-orthonormal-basis}。）
\end{enumerate}
\end{exercise}
% SUPPLEMENT: id=v14-ch05-ex62-solution; parent=ch05-ex62-lebesgue-l2-separable; kind=solution+motivation+hypotheses; origin=independent-derivation
\begin{solution}
\textbf{五问逐项补写（重构新增）。}
\begin{enumerate}
\item 先用\cref{theorem:ch02-simple-approximation}取简单函数 $s_k\to f$ 逐点且 $|s_k|\leqslant|f|$；支配收敛应用于 $|s_k-f|^2\leqslant4|f|^2$，得 $\|s_k-f\|_2\to0$。故只需逼近简单函数的每个示性函数。给 Lebesgue 可测 $E\subseteq[0,1]$，把 $\chi_E$ 零延拓到实线。由\cref{theorem:ch02-lone-density}，对任意 $\eta>0$，存在连续函数 $g$ 使 $\int_{\mathbb R}|g-\chi_E|<\eta^2$。把 $g$ 换成 $q=\min(1,\max(0,\operatorname{Re}g))$，它仍连续，取值在 $[0,1]$，且 $|q-\chi_E|\leqslant|g-\chi_E|$。因此
\begin{align*}
\|q-\chi_E\|_{L^2([0,1])}^2
\leqslant\int_{[0,1]}|q-\chi_E|<\eta^2.
\end{align*}
对有限和 $s=\sum_{j=1}^Na_j\chi_{E_j}$，逐项选连续 $q_j$，使 $\sum_j|a_j|\|q_j-\chi_{E_j}\|_2<\varepsilon/2$；例如每项误差小于 $\varepsilon/[2N(1+|a_j|)]$ 即可。再结合简单函数逼近，得到连续函数稠密。
\item 实多项式含常数、分离 $[0,1]$ 的点，由 Stone--Weierstrass 定理在实连续函数中一致稠密。分别逼近实部与虚部，得到复多项式在 $C([0,1])$ 中一致稠密。因区间测度为1，$\|g-p\|_2\leqslant\|g-p\|_u$；结合第一问得到 $L^2$ 稠密。
\item 系数属于 $\mathbb Q+\mathrm i\mathbb Q$ 的多项式只有可数多个。给 $p(t)=\sum_{j=0}^Na_jt^j$，以有理复数 $b_j$ 逐项逼近 $a_j$，则在 $[0,1]$ 上 $|p(t)-\sum b_jt^j|\leqslant\sum|a_j-b_j|$。故这可数集在全部多项式中一致稠密，再由第二问在 $L^2$ 中稠密。
\item 以 $[k,k+1)$（$k\in\mathbb Z$）分割实线。平移给每块 $L^2$ 与 $L^2([0,1))$ 等距同构；添去单个端点不改变 Lebesgue $L^2$，故每块可分。按一个自然数序列枚举这些块，应用上一道分解习题，得到 $L^2(\mathbb R)$ 可分。
\item 对 $n$ 归纳。$n=1$ 由第四问成立；若 $L^2(\mathbb R^{n-1})$ 可分，它和 $L^2(\mathbb R)$ 都有可数标准正交基，乘积基习题给 $L^2(\mathbb R^{n-1}\times\mathbb R)$ 的可数标准正交基，故可分。Euclid 空间 Lebesgue 测度是相应 Borel 乘积测度的完备化，完备化不改变 $L^2$ 等价类空间：由\cref{proposition:ch02-measurable-version-completion}，每个完备化可测函数都几乎处处等于原测度下的可测函数，积分相同。因此也可先把各因子等距识别为 Borel 测度下的 $L^2$，应用乘积基结论后再完备化，得到的正是 $L^2(\mathbb R^n)$。
\end{enumerate}
\end{solution}
\begin{remark}
\textbf{动机与逐条件说明（重构新增）。}逼近链为简单函数、连续函数、多项式、有理系数多项式，最终把不可数的函数空间压缩为一个可数稠密集。Lebesgue 测度用于连续逼近和端点零测；有限区间的有限测度把一致误差转为 $L^2$ 误差。第四问不能把全实线上的一致误差直接积分，而须先可数分块。第五问用有限维归纳和 $\sigma$ 有限的 Lebesgue 乘积测度，并区分乘积测度与其完备化。
\end{remark}
% END SUPPLEMENT: id=v14-ch05-ex62-solution
% SOURCE p=193; printed=178; object=ch05-ex63-weak-sphere; number=63; region=middle
\begin{exercise}\label{exercise:ch05-weak-sphere}
设 $\mathscr H$ 为无限维 Hilbert 空间。
\begin{enumerate}
\item\label{item:ch05-orthonormal-weak-zero} $\mathscr H$ 中每个标准正交序列都弱收敛于0。
\item\label{item:ch05-sphere-weakly-dense} 单位球面 $S=\{x:\|x\|=1\}$ 在单位闭球 $B=\{x:\|x\|\leqslant1\}$ 中弱稠密。（事实上，$B$ 中每个 $x$ 都是 $S$ 中某个序列的弱极限。）
\end{enumerate}
\end{exercise}
% SUPPLEMENT: id=v14-ch05-ex63-solution; parent=ch05-ex63-weak-sphere; kind=solution+motivation+hypotheses; origin=independent-derivation
\begin{solution}
\textbf{两问补写（重构新增）。}
\begin{enumerate}
\item 给标准正交序列 $(u_n)$ 和任意 $f\in\mathscr H^*$，Riesz 表示给 $f(x)=\langle x,y\rangle$。Bessel 不等式给 $\sum_n|\langle y,u_n\rangle|^2\leqslant\|y\|^2$，故各项趋零。由共轭对称性，$f(u_n)=\langle u_n,y\rangle\to0$。这正是弱收敛。
\item 固定 $x\in B$。其正交补 $\{x\}^\perp$ 无限维：若它有限维，由 $\mathscr H=\operatorname{span}\{x\}\oplus\{x\}^\perp$ 将推出 $\mathscr H$ 有限维。递归在该无限维空间选取标准正交序列 $(u_n)$；每一步先选不在已有有限张成空间中的向量，再正交化。令 $z_n=x+\sqrt{1-\|x\|^2}\,u_n$，则 Pythagoras 给 $\|z_n\|^2=1$。第一问给 $u_n\rightharpoonup0$，所以 $z_n\rightharpoonup x$。若 $\|x\|=1$，公式即常值序列。因而每个 $B$ 中向量都属于 $S$ 的弱闭包。由\subcref{exercise:ch05-weak-closures-completeness}{item:ch05-unit-ball-weakly-closed}，$B$ 弱闭且包含 $S$，故弱闭包恰为 $B$。
\end{enumerate}
\end{solution}
\begin{remark}
\textbf{动机与逐条件说明（重构新增）。}把缺少的长度放到不断变化的正交方向，任何固定连续线性测量都会看不到这部分误差。无限维用于持续选择新方向；有限维中弱拓扑与范数拓扑相同，球面闭，无法弱稠密于包含0的球。第一问不要求正交序列是基；第二问也不要求整个空间可分。范数始终为1却弱趋零，说明弱收敛与范数收敛不同。
\end{remark}
% END SUPPLEMENT: id=v14-ch05-ex63-solution
% SOURCE p=193; printed=178; object=ch05-ex64-shift-operators; number=64; region=middle-lower
\begin{exercise}\label{exercise:ch05-shift-operators}
设 $\mathscr H$ 为可分无限维 Hilbert 空间，标准正交基为 $\{u_n\}_{n=1}^{\infty}$。
\begin{enumerate}
\item\label{item:ch05-left-shift-strong-zero}
\begin{definition*}[左移算子]\phantomsection\label{definition:ch05-left-shift}
对 $k\in\mathbb N$，定义 $L_k\in L(\mathscr H,\mathscr H)$ 为
\begin{align*}
L_k\Bigl(\sum_{n=1}^{\infty}a_nu_n\Bigr)=\sum_{n=k+1}^{\infty}a_nu_{n-k}.
\end{align*}
\end{definition*}
则 $L_k\to0$ 在强算子拓扑中成立，但在范数拓扑中不成立。
\item\label{item:ch05-right-shift-weak-zero}
\begin{definition*}[右移算子]\phantomsection\label{definition:ch05-right-shift}
对 $k\in\mathbb N$，定义 $R_k\in L(\mathscr H,\mathscr H)$ 为
\begin{align*}
R_k\Bigl(\sum_{n=1}^{\infty}a_nu_n\Bigr)=\sum_{n=1}^{\infty}a_nu_{n+k}.
\end{align*}
\end{definition*}
则 $R_k\to0$ 在弱算子拓扑中成立，但在强算子拓扑中不成立。
\item\label{item:ch05-shift-products} $R_kL_k\to0$ 在强算子拓扑中成立，而对每个 $k$，$L_kR_k=I$。（使用\subcref{exercise:ch05-strong-operator-products}{item:ch05-strong-product-convergence}。）
\end{enumerate}
\end{exercise}
% SUPPLEMENT: id=v14-ch05-ex64-solution; parent=ch05-ex64-shift-operators; kind=solution+motivation+hypotheses; origin=independent-derivation
\begin{solution}
\textbf{三种收敛的完整比较（重构新增）。}正交展开唯一，且 $(a_n)\in\ell^2$，所以题中的两个级数均范数收敛，定义良好。坐标线性给两个算子线性。
\begin{enumerate}
\item 对 $x=\sum a_nu_n$，$\|L_kx\|^2=\sum_{n>k}|a_n|^2\leqslant\|x\|^2$，且尾和趋零，故 $L_k$ 有界并强趋零。但 $L_ku_{k+1}=u_1$，故 $\|L_k\|=1$，不按算子范数趋零。
\item $\|R_kx\|^2=\sum_n|a_n|^2=\|x\|^2$，所以 $R_k$ 为等距算子。若 $y=\sum b_nu_n$，
\begin{align*}
|\langle R_kx,y\rangle|
=\Bigl|\sum_{n=1}^{\infty}a_n\overline{b_{n+k}}\Bigr|
\leqslant\|x\|\Bigl(\sum_{j>k}|b_j|^2\Bigr)^{1/2}\longrightarrow0.
\end{align*}
该标量级数由 Schwarz 绝对收敛；Riesz 表示使所有连续线性泛函都被此测试覆盖，故 $R_k\to0$ 于弱算子拓扑。对任意非零 $x$，$\|R_kx\|=\|x\|$，所以不强趋零。
\item $\|R_kL_kx\|=\|L_kx\|\to0$，故第一结论成立。直接由坐标公式，右移后再左移 $k$ 位恢复每个系数，故 $L_kR_kx=x$ 对所有 $x$。原提示也可通过\subcref{exercise:ch05-strong-operator-products}{item:ch05-strong-moving-input}的估计思想来理解：这里输入 $L_kx\to0$，外层算子满足统一界 $\|R_k\|=1$，所以像仍趋零；并不需要 $R_k$ 自身强收敛。
\end{enumerate}
\end{solution}
\begin{remark}
\textbf{提示与条件说明（重构新增）。}原提示可以用于第二问的另一种反证：若 $R_k$ 也强趋零，因 $L_k$ 已强趋零，前题乘积结论就会给 $L_kR_k\to0$ 强收敛，这与 $L_kR_k=I$ 及空间非零矛盾。证明 $R_kL_k\to0$ 时则直接使用统一有界估计，不能把尚不成立的 $R_k$ 强收敛当成已知条件。可分性和无限维保证基可编号为无限序列，标准正交性把范数变成平方和；完整性保证每个输入都有唯一展开。这个例子同时区分算子范数、强算子和弱算子收敛，并说明只有弱算子收敛时交换乘法和极限会失败。
\end{remark}
% END SUPPLEMENT: id=v14-ch05-ex64-solution
% SOURCE p=193; printed=178; object=ch05-ex65-hilbert-dimension; number=65; region=lower
\begin{exercise}\label{exercise:ch05-hilbert-dimension}
$\ell^2(A)$ 与 $\ell^2(B)$ 酉同构，当且仅当 $\card(A)=\card(B)$。
\end{exercise}
% SUPPLEMENT: id=v14-ch05-ex65-solution; parent=ch05-ex65-hilbert-dimension; kind=solution+motivation+hypotheses; origin=independent-derivation
\begin{solution}
\textbf{任意基数情形补写（重构新增）。}若存在双射 $\phi:A\to B$，令 $(Uf)(\beta)=f(\phi^{-1}(\beta))$。它线性满射，非负平方和在指标重排下不变，故等距，由极化判别为酉映射。

反之，设 $U:\ell^2(A)\to\ell^2(B)$ 酉。令 $v_\alpha=Ue_\alpha$；保持内积使其标准正交，且若 $y$ 与每个 $v_\alpha$ 正交，写 $y=Ux$ 得 $x$ 与每个 $e_\alpha$ 正交，于是 $x=0$，所以 $\{v_\alpha\}_{\alpha\in A}$ 是 $\ell^2(B)$ 的标准正交基。

若一个指标集有限，其空间就是相应有限维坐标空间。可逆线性映射保持有限维维数：独立组的像独立，基的像仍张成。若另一指标集无限，它的标准单位向量含任意大的有限独立组，矛盾；故两个都有限且基数相等。这也涵盖允许空指标集时的零空间。

现在设 $A,B$ 都无限。对每个 $\beta\in B$，Bessel 不等式给 $A_\beta=\{\alpha:\langle e_\beta,v_\alpha\rangle\ne0\}$ 至多可数。每个非零 $v_\alpha$ 至少有一个非零坐标，故 $A=\bigcup_{\beta\in B}A_\beta$，从而 $\card(A)\leqslant\card(B)\cdot\aleph_0=\card(B)$。交换两个基的角色：每个 $\alpha$ 对应的 $B_\alpha=\{\beta:\langle v_\alpha,e_\beta\rangle\ne0\}$ 至多可数，而 $\{v_\alpha\}$ 完全使 $B=\bigcup_\alpha B_\alpha$，故 $\card(B)\leqslant\card(A)$。由 Cantor--Bernstein 定理得两者相等。
\end{solution}
\begin{remark}
\textbf{动机与逐条件说明（重构新增）。}Hilbert 空间的酉类型由标准正交方向的数目决定，不能只用代数基维数替代。无限基数比较依赖“每个向量的非零正交坐标至多可数”，而有限情形必须单独处理，因为 $n\cdot\aleph_0$ 无法恢复有限 $n$。这里沿用本书选择原理及无限基数乘法的约定，不要求 $A,B$ 可数。保持复线性、内积及满射使基完整地对应；单纯等距嵌入只会给单向包含信息。
\end{remark}
% END SUPPLEMENT: id=v14-ch05-ex65-solution
% SOURCE p=193; printed=178; object=ch05-ex66-closed-continuous-l2-subspace; number=66; region=lower
\begin{exercise}\label{exercise:ch05-closed-continuous-l2-subspace}
设 $\mathscr M$ 为 $L^2([0,1],m)$ 的闭子空间，且 $\mathscr M\subseteq C([0,1])$。
\begin{enumerate}
\item\label{item:ch05-l2-controls-uniform} 存在 $C>0$，使 $\|f\|_u\leqslant C\|f\|_{L^2}$ 对所有 $f\in\mathscr M$ 成立。（使用闭图像定理。）
\item\label{item:ch05-continuous-evaluation-representers} 对每个 $x\in[0,1]$，存在 $g_x\in\mathscr M$，使 $f(x)=\langle f,g_x\rangle$ 对所有 $f\in\mathscr M$ 成立，且 $\|g_x\|_{L^2}\leqslant C$。
\item\label{item:ch05-continuous-subspace-dimension} $\mathscr M$ 的维数至多为 $C^2$。（提示：若 $\{f_j\}$ 为 $\mathscr M$ 中的标准正交序列，则对所有 $x\in[0,1]$，$\sum_j|f_j(x)|^2\leqslant C^2$。）
\end{enumerate}
\end{exercise}
% SUPPLEMENT: id=v14-ch05-ex66-solution; parent=ch05-ex66-closed-continuous-l2-subspace; kind=solution+motivation+hypotheses; origin=independent-derivation
\begin{solution}
\textbf{三问完整补写（重构新增）。}先说明题设的识别：每个 $\mathscr M$ 中的 $L^2$ 等价类有一个连续代表，并且唯一。若两个连续函数几乎处处相同而在某点不同，连续性给一个相对开区间，在其上差的模严格为正；该区间有正 Lebesgue 测度，矛盾。
\begin{enumerate}
\item 令 $J:\mathscr M\to C([0,1])$ 把等价类映为唯一连续代表，显然线性。两端均 Banach：前者是完备 $L^2$ 的闭子空间，后者取一致范数。若 $f_n\to f$ 在 $\mathscr M$ 的 $L^2$ 范数下，且 $Jf_n\to g$ 一致，则因 $m([0,1])=1$，也有 $Jf_n\to g$ 于 $L^2$。极限唯一给 $f=[g]$，连续代表唯一给 $Jf=g$。因此 $J$ 的图像序列闭；乘积是度量空间，故图像闭。由\cref{theorem:ch05-closed-graph}，$J$ 有界，取任意正数 $C\geqslant\|J\|$ 即得结论。
\item 固定 $x$，评价泛函 $E_x(f)=f(x)$ 线性且 $|E_x(f)|\leqslant\|f\|_u\leqslant C\|f\|_2$，所以其范数至多为 $C$。在 Hilbert 空间 $\mathscr M$ 上用 Riesz 表示定理，得唯一 $g_x\in\mathscr M$，满足所述表示和 $\|g_x\|_2=\|E_x\|\leqslant C$。
\item 取 $\mathscr M$ 中任意有限标准正交族 $f_1,\ldots,f_N$。对每个 $x$，用第二问及 Bessel 不等式，
\begin{align*}
\sum_{j=1}^N|f_j(x)|^2
=\sum_{j=1}^N|\langle f_j,g_x\rangle|^2
=\sum_{j=1}^N|\langle g_x,f_j\rangle|^2
\leqslant\|g_x\|_2^2\leqslant C^2.
\end{align*}
在 $[0,1]$ 上积分，得到 $N=\sum_{j=1}^N\|f_j\|_2^2\leqslant C^2$。若空间维数大于 $C^2$，可取整数 $N>C^2$ 个线性无关向量，再有限 Gram--Schmidt 得到这样大的标准正交族，矛盾。故空间有限维且维数至多为 $C^2$。对无限标准正交序列的原提示，由上述有限和界取上确界即可。
\end{enumerate}
\end{solution}
\begin{remark}
\textbf{动机与逐条件说明（重构新增）。}闭图像定理先把“每个元素都有连续代表”提升为统一范数控制，再把点评价表示为空间内的向量，最后积分 Bessel 界限制独立方向数。$L^2$ 闭性负责 $\mathscr M$ 完备；若取所有多项式组成的非闭子空间，则维数无限，结论失败。连续代表使点评价有定义，Lebesgue 测度的满支撑使该代表唯一；一般 $L^2$ 类不能随意在一个点求值。区间测度为1使维数界恰为 $C^2$；全体 $x$ 使用同一个常数至关重要。允许零空间，不必假设已有无限正交序列。
\end{remark}
% END SUPPLEMENT: id=v14-ch05-ex66-solution
% SOURCE p=193; printed=178; object=ch05-ex67-mean-ergodic; number=67; region=bottom
\begin{exercise}[平均遍历定理]\label{exercise:ch05-mean-ergodic}
设 $U$ 为 Hilbert 空间 $\mathscr H$ 上的酉算子，$\mathscr M=\{x:Ux=x\}$，$P$ 为到 $\mathscr M$ 的正交投影（\cref{exercise:ch05-orthogonal-projection}），且 $S_n=n^{-1}\sum_{j=0}^{n-1}U^j$。则 $S_n\to P$ 于强算子拓扑。（若 $x\in\mathscr M$，则 $S_nx=x$；若对某个 $y$ 有 $x=y-Uy$，则 $S_nx\to0$。由\subcref{exercise:ch05-hilbert-adjoint}{item:ch05-unitary-adjoint-inverse}，$\mathscr M=\{x:U^*x=x\}$。对 $T=I-U$ 应用\subcref{exercise:ch05-hilbert-adjoint}{item:ch05-hilbert-adjoint-range-kernel}。）
\end{exercise}
% SUPPLEMENT: id=v14-ch05-ex67-solution; parent=ch05-ex67-mean-ergodic; kind=solution+motivation+hypotheses; origin=independent-derivation
\begin{solution}
\textbf{完整补写（重构新增）。}$\mathscr M=\ker(I-U)$ 为闭子空间，所以 $P$ 确存在。因为 $U^*=U^{-1}$，有 $Ux=x$ 当且仅当 $U^*x=x$。于是伴随的值域与核关系给
\begin{align*}
\mathscr R(I-U)^\perp=\ker((I-U)^*)=\ker(I-U^*)=\mathscr M,
\qquad \overline{\mathscr R(I-U)}=\mathscr M^\perp,
\end{align*}
后一等式再用双正交补习题。对 $x\in\mathscr M$，每个 $U^jx=x$，故 $S_nx=x=Px$。若 $x=(I-U)y$，则望远镜求和给
\begin{align*}
S_nx=\frac1n\sum_{j=0}^{n-1}(U^jy-U^{j+1}y)
=\frac{y-U^ny}{n},\qquad
\|S_nx\|\leqslant\frac{2\|y\|}{n}\longrightarrow0.
\end{align*}
此外 $\|S_n\|\leqslant n^{-1}\sum_{j=0}^{n-1}\|U^j\|\leqslant1$。对任意 $z\in\mathscr M^\perp$ 和 $\varepsilon>0$，选 $w\in\mathscr R(I-U)$ 使 $\|z-w\|<\varepsilon/2$；再选 $n$ 足够大使 $\|S_nw\|<\varepsilon/2$，于是 $\|S_nz\|\leqslant\|z-w\|+\|S_nw\|<\varepsilon$。故 $S_nz\to0$ 于整个 $\mathscr M^\perp$。最后对任意 $x=Px+(x-Px)$，有 $S_nx=Px+S_n(x-Px)\to Px$，即强算子收敛。
\end{solution}
\begin{remark}
\textbf{动机与逐条件说明（重构新增）。}平均保留每次迭代都不变的分量；形如 $y-Uy$ 的差分在求和中只剩两个端点，除以 $n$ 后消失。酉性同时给迭代的一致范数界及 $U^*=U^{-1}$，从而准确识别差分值域闭包。Hilbert 结构和闭性负责正交分解；不能只在值域上证收敛便跳到其闭包，统一界补上了这一步。结论对每个固定向量成立，并未声称算子范数收敛，也未要求空间可分。零空间与 $U=I$ 的情形都包含在证明中。
\end{remark}
% END SUPPLEMENT: id=v14-ch05-ex67-solution
\end{document}
